\begin{afterword}
\summaryhead

The post starts from the author's habit, while writing the Albert and Barry dialogue, of saying ``Either the falling tree makes a sound, or it does not!'' It argues that ``P or not-P'' is not a safe rule for English sentences: the liar sentence and ``Have you stopped beating your wife?'' fail it, and a classical logician can keep the rule only by subtle moves.

The rule needs P to mean exactly the same thing in both halves. If Albert's ``sound'' differs from Barry's, the tree can make one and not the other, which the author writes as Albert::sound and Barry::sound, borrowing C++ notation. Many words are functions of hidden context variables: there is no ``the left side'' of a street, only your left as you travel. When a hidden variable takes different values, reality seems unstable, which the author says confuses undergraduates and postmodernist professors. Strictly, a sentence with a free variable, like ``2 + 2 = X,'' has no truth value at all.

\responsehead

The post's main point is correct and useful: before asking whether a statement is true, check that it means the same thing each time it is used. The grocery-store example shows it well, and the closing correction, that ``2 + 2 = X'' has no truth value rather than a variable one, is the standard logician's point, stated accurately.

The point is also old, and the post does not mention its old name. Using a term in one sense at one place and another sense at another is the fallacy of equivocation, which Aristotle listed among the fallacies that depend on language. The post coins a new name for it. What it adds is the extension to context-dependent words such as ``left,'' where the shifting meaning comes from a variable no one states.

The claim that a shifting variable makes reality itself feel unstable rests on the account of categorization in ``How An Algorithm Feels From Inside,'' and inherits its lack of evidence about brains.

Finally, the post assigns the confusion to ``undergraduates (and postmodernist professors)'' without naming or quoting one. Its only example of their reasoning is a speech the author composed, and the analysis of it is ``left as an exercise to the reader.''

In short: an old fallacy under a new name, correctly stated and usefully extended, with an opposing view represented only by a speech the author wrote for it.
\end{afterword}
